\vspace{-2mm}
\section{Introduction}
\vspace{-2mm}
% Recent progress has demonstrated impressive capabilities in simulation and on hardware, but current approaches still suffer from three major limitations. First, they remain task-specific: most policies are trained to imitate motion capture clips or solve a single locomotion or manipulation task. Second, they are non-adaptive: once trained, policies cannot be easily fine-tuned or composed for new tasks. Third, they lack a unified and explainable interface for goal specification and behavior composition, making it difficult for human operators to direct the robot or combine learned skills into new behaviors.

% To overcome these limitations, we take inspiration from the success of large language models (LLMs) and vision-language models (VLMs), which pair large-scale pretraining with post-training alignment and test-time adaptation to create versatile general-purpose systems.

% To overcome these limitations, we take inspiration from the success of large language models (LLMs) and vision-language models (VLMs), which pair large-scale pretraining with post-training alignment and test-time adaptation to create versatile general-purpose systems.

% Humanoid is important and WBC is important
Humanoid robots have the potential to transform numerous aspects of our daily lives, from manufacturing and logistics to healthcare and personal assistance. However, realizing this potential requires robots to perform a wide range of tasks in dynamic and unstructured environments. Humanoid whole-body control is a fundamental and challenging problem in robotics, serving as the first step to enable the humanoids to work safely in human environments~\citep{gu2025humanoid}.

% Robotics foundation models
In robotics, foundation models have the potential to unify diverse control objectives under a single policy, allowing robots to adapt to new tasks in a zero-shot\footnote{\emph{Zero-shot} means that, after pre-training, the policy can be directly deployed in the real world without further interacting with either simulated or real environments. In contrast, \emph{few-shot} means the policy needs to interact with the environment to collect new data in few episodes to improve on certain tasks.}way or with efficient post-training. The closest approaches to such paradigms are Vision-Language-Action (VLA) models for robotic manipulations~\citep[e.g.,][]{Ghosh2024octo,pi05,openvla,Zhong2025DexGraspVLA, geminirobotics,grootn1} that learn from human demonstrations (i.e., behavior cloning). However, for humanoid whole-body control, there is a fundamental mismatch that limits direct behavior cloning: unlike manipulation tasks, there are no readily available actuator-level action labels or large-scale teleoperation datasets.
 
 
 % , as the networks are typically optimized for manipulation, where control frequencies can remain relatively low (e.g., issuing a grasp or push command at 5–10 Hz) without compromising task success. By contrast, whole-body humanoid control demands high-frequency actuation (often 50-200 Hz) to maintain balance. 
 
For whole-body humanoid control, most recent advancements follow the sim-to-real pipeline and rely on reinforcement learning (RL) to train policies in simulation before transferring them to hardware~\citep{gu2025humanoid}. Following the success of RL-based motion tracking in physics-based character animation~\citep[e.g.,][]{Luo2024universal, TesslerGNCP24,TirinzoniTFGKXL25zeroshot}, recent works~\citep[e.g.,][]{Zakka2025mujocoplayground,Seo2025fasttd3,Chen2025gmt,liao2025beyondmimic,he2025asap,cheng2024expressive,he2025omnih2o} have shown remarkable results in transferring policies trained in simulation to real robots. However, most of these approaches rely on \emph{on-policy policy gradient} methods (e.g., PPO~\citep{schulman2017proximal}) with \emph{explicit tracking-based rewards} and suffer from major limitations. First, they remain task-specific: most policies are trained to explicitly imitate motion capture clips or solve a single task. Second, they are non-adaptive: once trained, policies cannot be easily fine-tuned or composed for new tasks. Third, they lack a unified and explainable interface for goal specification and behavior composition, making it difficult for human operators to direct the robot or combine learned skills into new behaviors.

%While almost the totality of existing approaches for real-world deployment leverage on-policy training, we 

In this work, we investigate whether \emph{off-policy unsupervised} RL can be a suitable approach to train so-called Behavioral Foundation Models (BFMs) for whole-body control of a humanoid robot, enabling it to solve a wide range of downstream tasks specified by rewards, goals, or demonstrations without retraining. For tasks that require retraining, the BFM should enable efficient post-training. 
% \guanya{todos: add one sentence for efficient post-training}
This conjecture is far from trivial. First, most existing methods with real-world deployment rely on on-policy training (primarily PPO), and there is little evidence that off-policy learning—commonly used in unsupervised RL for training multi-task policies—is well suited to this context. Second, no evidence exists that unsupervised RL algorithms can handle the sim-to-real gap and dynamic disturbances robustly, either during simulation policy training or at real-world inference.

% In particular, we build on the FB-CPR algorithm recently proposed by~\citet{TirinzoniTFGKXL25zeroshot} to control a virtual physics-based humanoid character and we investigate how it can be used to train a BFM for whole-body control of a humanoid robot. 
We develop \method\footnote{\texttt{\textcolor{myred}{\textbf{Zero}}} comes from its zero-shot inference capability via unsupervised RL and it is a first-of-its-kind model.}, an online off-policy unsupervised RL algorithm that leverages motion capture data to regularize the process of learning generalist whole-body control policies towards \emph{human behaviors}. 
We introduce domain randomization to address the sim-to-real gap and train robust policies via asymmetric history-dependent training, leveraging the privileged information available in simulation. Additionally, we incorporate auxiliary rewards to ensure that the learned behaviors adhere to the safety and operational constraints of the physical robot. To the best of our knowledge, the resulting algorithm allows us to train the \emph{first behavioral foundation model} for real humanoids that can be prompted for different tasks (e.g., reward optimization, pose reaching, and motion tracking) without retraining (i.e., in zero-shot). Such a flexible and ready-to-use model, paves the way to fast adaptation, fine-tuning or even high-level planning.
%
We validate our approach in both simulated environments and on a real Unitree G1 humanoid (Fig.~\ref{figure_1_teaser} for examples), demonstrating robust generalization across tasks and conditions, and showing that even when the zero-shot policy is not satisfactory, we can effectively improve it within a few episodes of environment interaction. The discussion of related work is available in \Cref{sec:appendix-related-work}.


% OLD
% In this work, we look into unsupervised reinforcement learning (RL) as a promising paradigm to train so-called \emph{behavioral foundation models} to enable agents to solve a wide range of downstream tasks specified by rewards, goals, or demonstrations without any retraining. In particular, we build on the FB-CPR algorithm recently proposed by~\citet{TirinzoniTFGKXL25zeroshot} to control a virtual physics-based humanoid character and we investigate how it can be used to train a BFM for whole-body control of a humanoid robot. 
% %trained in simulation and transferred zero-shot to a physical robot. Our model is designed to represent behaviors in a latent space shared across states, goals, and rewards, enabling a single policy to be prompted for various tasks (such as motion tracking, goal-reaching, and reward optimization) without retraining. We build on top of the unsupervised RL approach of~\citet{TirinzoniTFGKXL25zeroshot} and address the challenges of sim-to-real transfer. 
% We introduce domain randomization to address the sim-to-real gap and train robust policies via asymmetric history-dependent training to leverage the privileged information available in simulation. Additionally, we incorporate reward regularization to ensure that the learned behaviors adhere to safety and operational constraints of the physical robot. The resulting algorithm allows us to train the \textbf{first behavioral foundation model for the Unitree G1} that is light --the policy can be deployed on the robot's onboard computer-- and general --it can be prompted for different tasks (e.g., reward optimization, pose reaching, and motion tracking) without retraining (i.e., in zero-shot). Such a flexible and ready-to-use model, paves the way to fast adaptation, fine-tuning or even high-level planning, thereby advancing the development of versatile humanoid robots with general-purpose capabilities. 
% %
% We validate our approach in both simulated environments and on a real Unitree G1 robot, demonstrating robust generalization across tasks and conditions, and showing that even when the zero-shot policy is not satisfactory, we can effectively improve it in a matter of few simulated or real episodes.